\section{The Labors of Hercules, Part 3}

\paragraph{The Third Test, The Ceryneian Hind}

For the third labor\footnote{\url{https://en.wikipedia.org/wiki/Ceryneian\_Hind}}, Hercules was given a retrieval task instead of a slaying to accomplish. Since Hercules could not be overcome with guile
\& brute force, it was hoped that he could be made to trespass against a god, \& have divine fury invoked upon him. Specifically, Eurystheus hoped that this labor would infuriate Artemis, who would presumably give him the terrible fate that Pentheus\footnote{\url{https://en.wikipedia.org/wiki/Pentheus}} met with on the Holy Mountain.The hind, additionally, traveled faster than an arrow. How was it to be supposed that Hercules could conquer both the speed of flight \& the jealousy of the female Huntress?

Since deer did not inhabit Greece, this myth carries an echo of Northern Lands. The white or golden stag is not merely a supernatural symbol, but a regal one. Hercules sees the deers antlers' glinting very far off, \& begins the chase, which lasts (some say) an entire year, possibly in the Hyperborean north. At last, either by a ruse or skill of the arrow (possibly shooting one between its legs to trip it, or by using a net) or by direct permission of Artemis herself (who sides with the rugged and manly hero), Hercules comes into possession of the golden deer, \emph{on the promise that he will return it to Artemis when he is done}.

When he brings it to Eurystheus, who intends to keep the deer, Hercules wisely lets the deer go seconds before delivering it up -- the deer slips away, \& the mighty Herc shrugs it off: ``You aren't fast enough.''

Hercules proves himself a man of honor in this task, as he does not slay or injure the deer, but returns it as promised to its rightful owner. Here we should note that, in esoteric action, as in any other portion of Life, there are laws that obtain. Although those who travel in the astral realm encounter far fewer laws than we do, there are strictures that (for all that) may be even more indurate than gravity. One of these, according to \emph{Gnosis} (Mouravieff) is that a man cannot advance in practice unless he trains or leaves behind a replacement. Karma is another law that obtains; repentance can mitigate or erase it, but still, there is something that holds true \& fast -- in this case, another bears your sins. Or who did you think paid your trespass?

I say this to encourage readers to not set up a dichotomy in their mind between royal liberation and the lesser mysteries. Not until one is set free from Time \& Space itself does one become ``free of all Law'', not even if one is ``awakened'' fully into the etheric, astral, and spiritual realms. Besides, the dichotomy itself is deadening. If one is truly free, is one not silent? Hercules is careful to observe piety and honor; it is not forced, out of fear: he speaks and converses with Artemis almost as an equal, as a better, but someone he may speak before, and she grants him favor in her eyes. There is one equal to, and worthy of, capturing the devoted and precious hind. It is the pious warrior, who nevertheless speaks with ``Frankness'' before the very gods. St. Paul spoke of this when he said, ``all things are lawful, not all things are useful''.

Hercules is wise enough to insult and belittle his enemy; after all, if one doesn't add a little salt to the wound, who is going to arrange the rest of the tests? Such a jibe befits the warrior. It is the jibe of the Russian general Kutuzov to the French prisoners of Napoleon's army, after he has magnanimously spoken of forgiveness. Rallying his own troops, he says to them: ``But after all, who asked them to come here, anyway?''. It is Frederick the Great shouting out to his fleeing men, ``Ihr wolt, ewige leben?!'' It is Brennus before the conquered Romans, speaking very simply back to them something that would become their own watchword- \emph{Vae Victis}. Forgiveness and detachment are the rocks on which the raging sea breaks, and falls, dashing those who ride it. There should be a kind of sacrosanct danger about sinning against a truly pious man; when defers his own judgement and wrath, the wicked should tremble. Archangel Michael, no doubt in a very great wroth and sorely tempted to pass his bounds, declared to Satan's impudence, ``May the Lord judge between us!''. The wrath of the righteous, the fury of the hero\footnote{\url{https://www.youtube.com/watch?v=ggFKLxAQBbc}}, the anger of the good man\footnote{\url{https://www.youtube.com/watch?v=9ClCMR58LV8}} pushed too far, should resolve into that ritual and controlled resistance, a detached willingness to see it through bitterly, without bitterness, which heaps coals of fire upon the head of those who take the part of Satan. This is how the enigma of forgiveness \& resistance is resolved -- in the action of the hero, who must understand both. Could he not forgive, he could not formalize his actions and discipline them to undergo a trial, but would simply ``rage out'' and go for the throat of his persecutor. Could he not resist, he would forever remain under their feet. This tension drives the warrior, like a bow drives the arrow. To those who cannot see the union of this, they are either not warriors, or do not see that ``the insanity of God is greater than the wisdom of man'' (Plato).

Guided by the wisdom of a Solomon, forgiveness and power are inseparable. Are not detachment and passion reconciled by the warrior? When someone loses their temper and lashes out in anger and injustice against you, the ability to bear the stroke without retributive ire in the same manner is actually an esoteric technique for transferring the energy they are losing to yourself. This is the secret meaning of ``heaping coals of fire'' upon their head -- a person who gives in to the ``wrath of man'' loses enormous reserves of spiritual energy : especially if some physical scapegoating occurs, those energies often go to the victim, rightfully. Received in the right frame of mind, they become available as a ``lost talent''. This is why God invites the ``poor'' of the world to his banquet : the rich folk turned him down. Their loss, our gain. If Christianity is the religion of the Kali Yuga, do not therefore conceive that it is bound to share its inadequacies: it is made to triumph over the greatest enemy, the last enemy, to devour it inside out.

If it wasn't the Kali Yuga, we wouldn't be here, reading and talking to each other to determine the fate of the new world. Other, better men would be leading. ``The better man'' in the end is just the one who does more things that are difficult for him -- if that is the test, it is a blessed time to be alive. He who plants a garden in the Kali Yuga, against all odds, stands above he who conquers worlds, with some odds in his favor.

``The last shall be first --- greater are those who have not seen, yet believed.''

Hercules triumphs by his ruse, his skill, his polity, his innate dignity as a true man, before man, beast, or gods. He is not to be deterred or thwarted. Nor will he stoop, except to make a well-earned jibe to further incite his opponent to useful wrath.

He doesn't show up off the street and demand to be Hercules. He IS Hercules, by virtue and dint of a long process, created in heaven, ratified in the mud of earth. He has conquered mental fog and the physical passions. He has risen above those things which men almost never even begin to subdue. He is already heads and shoulders above all men, and even though, as a god, he could use his liberty and peerage, he does not. ``Not counting himself to be equal to God, he humbled himself... '' Hercules prefigures the gracious and valiant and terrible true knight Jesus, who is far from the lamentable and tragic figure some neo-pagans think him to be. Hercules is the new Sun\footnote{\url{https://www.youtube.com/watch?v=gEaa9oAwT5U}}. So, in this third test which seems not so deadly because it was so successful, Hercules wins the affection and protection of Artemis; he is effectively adopted. The gods themselves are beginning to take sides. Heaven is being moved by earth. Hercules has swayed an eternal-feminine power to his side, and a mighty one at that, Artemis-Diana of the Hunt. It is a wild and primeval power that is now declared for the hero, who has presumed nothing, but has simply stood as a man ought to stand. As Hercules' divinity grows more obvious, the divinities who favor him, and yet remain worshiped and honored by him, begin to light up. It is as if they are lights which come alive and shine upon the hero, cutting his contours out of the dark, healing and supporting him. He grows greater, as he makes the gods greater than he. Further, he grows greater than the gods. In the eternal moment at the end of time, when Christ delivers the kingdom of God back to the hands of the Father, the eternal warrior announces the end of all things and annihilates the worlds with the breath of the Father. This is how the warrior becomes God: with the gesture of the true man, who yields fealty, forever indominant, to the rightful and original Lord. This makes Him the destroyer of worlds, for it is this act which is the deeper magic. Our world does not understand or accept this: that is why ``the wisdom of the world is foolishness with God''. If Evil cannot stand in the days of mercy, how shall it stand in the hour of judgement, which comes? What will Eurystheus do? He is running out of labors. Hercules is proving a difficult adversary. Satan himself may have to enter the lists, as there is a ``man worth killing''\footnote{\url{https://www.youtube.com/watch?v=ELoFRrUrcYQ}}.

Illustration:

Berlin, Neues Museum Herkules besiegt die goldbekrönte Hirschkuh (Herkules fängt die Hirschkuh von Ceryneia) Maler: Adolf Schmidt 

\flrightit{Posted on 2013-08-24 by Logres}

\begin{center}* * *\end{center}

\begin{footnotesize}
\begin{sffamily}
\texttt{Michel on 2013-08-25 at 12:38 said:}

``We swim against the rising waves, and crash against the shore, your body bends until it breaks, the early morning sings no more.

So rest your head its time to sleep, and dream of what's in store, your body bends until it breaks, and sings again no more, coz time has torn the flesh away, the early morning sings no more.''


\hfill

\end{sffamily}
\end{footnotesize}

